\subsection{Strong mixing at the physical scale and the averaged
condition}\label{rW9:ssec}

For spin systems with finite-range interactions, a spatial mixing condition
that holds uniformly in the boundary condition gives a spectral gap of the
Glauber dynamics that is uniform in the volume. Martinelli and Olivieri
derive it from strong mixing on large cubes~\cite{MO94a,MO94b}, Stroock and
Zegarlinski from the Dobrushin--Shlosman mixing condition~\cite{SZ92a,SZ92b},
and Cesi and Bertini, Cancrini and Cesi through quasi-factorisation of the
variance and the entropy~\cite{Ces01,BCC02}; see~\cite{Mar99,GZ03} for
surveys. The averaged condition $\mathrm{BP}(n,\varpi)$ is a statement of this
kind for the dynamics that resamples whole cubes of side $n$. For a block
dynamics with blocks larger than the mixing length, Dyer, Sinclair, Vigoda
and Weitz derive the contraction of the block dynamics from strong spatial
mixing by path coupling~\cite{DSVW04,BD97}, in the spirit of the
constructive criterion of~\cite{DS85}. This subsection carries that
argument out for the cube chain of Sec.~\ref{ssec:averaged} on gauge orbits.
It formulates a mixing condition on the cubes of side $n$ alone, posed for
couplings of the cube laws modulo gauge transformations and counted in cells
of side $s$, and proves that it implies $\mathrm{BP}(n,\varpi)$ with $\varpi$
bounded below by an explicit function of the mixing constants and of $s/n$;
with $n=\ell/a$ and $s=\ell_c/a$ the bound does not depend on $a$
(Theorem~\ref{rW9:main}, Corollary~\ref{rW9:physical}). The spectral step is
that of Corollary~\ref{cor:bpcoupling}; what is new is that the contraction
is required only for configurations that differ in one cell, so that the
hypothesis is posed on single cubes. The subsection then shows where the
condition holds and where it fails, and how it relates to (D) and to
$\mathrm{BP}$. The condition is not proved along the continuum tuning for any
group.

\emph{Setting.} $G$, $\delta$, $\varphi$, $\Lambda$, $E$, $P$, $\mu$, $V_x$,
$X$, $\partial V$, $\gamma^0_V$, $\mathcal P_x$, $\mathcal A$ and $\mathcal K$
are as in Sec.~\ref{ssec:averaged}; for $G=\SU(N)$ on a torus, $Z^\nu_t$ and
$\mathcal A_0$ are as in Sec.~\ref{ssec:bpweak}. Write $\mathcal G$ for the
gauge group, $G^{\Lambda}$ acting on $G^E$ by $(gU)_{(y,\nu)}=g_yU_{(y,\nu)}
g_{y+e_\nu}^{-1}$, and $\mathcal Z$ for the finite abelian group generated by
the maps $Z^\nu_t$. For $x\in X$ let $\Gamma_x(U,\cdot)$ be the law on $G^E$
of the configuration equal to a sample of $\gamma^0_{V_x}(\cdot\mid U)$ on
$V_x$ and to $U$ off $V_x$, the transition kernel of the move at $x$; so
$\mathcal P_xf(U)=\int f\,d\Gamma_x(U,\cdot)$ and $\mathcal K=\lvert X\rvert
^{-1}\sum_x\Gamma_x$. Every $\gamma^0_V(\cdot\mid U)$ has a continuous
positive density with respect to the Haar measure of $G^V$, and so have its
marginals.

\emph{Cells.} Fix an integer $s\ge1$; on $\Lambda_L$ assume that $s$ divides
$L$. For $c\in s\Z^4$ (read modulo $L$ on the torus) let $Q_c$ be the set of
links of $E$ with base point in $c+[0,s)^4$. The nonempty sets $Q_c$ are the
\emph{cells}; they partition $E$, and $\mathcal Q$ denotes their set. The
\emph{grid distance} of $Q_c$ and $Q_{c'}$ is $\lvert Q_c-Q_{c'}\rvert_s=
\lvert c-c'\rvert_\infty/s$, with differences on $\Lambda_L$ taken in
$(-L/2,L/2]^4$. The grid points at $\ell^\infty$-distance exactly $js$ from
$c$ number $(2j+1)^4-(2j-1)^4=64j^3+16j$ for $j\ge1$, so at most $64j^3+16j$
cells lie at grid distance $j$ from a given cell. A cell $Q$ and a point $x\in
X$ \emph{touch} if $Q\not\subset V_x$ and $Q\cap(V_x\cup\partial V_x)\ne
\emptyset$.

For $U,U'\in G^E$ let $\mathrm{H}_s(U,U')$ be the number of cells $Q$ with
$U_Q\ne U'_Q$, and
\[
\begin{split}
d_{\mathcal G}(U,U')&=\min_{g\in\mathcal G}\mathrm H_s(U,gU'),\\
d_{\mathcal{GZ}}(U,U')&=\min_{g\in\mathcal G,\,Z\in\mathcal Z}\mathrm H_s(U,
gZU'),
\end{split}
\]
the second for $G=\SU(N)$ on a torus. The minima are attained, since $\mathrm
H_s$ takes integer values. For a pseudometric $d$ on $G^E$ and probability
measures $\alpha,\beta$ write $W_1^d(\alpha,\beta)=\inf_\pi\int d\,d\pi$, the
infimum over couplings $\pi$ of $\alpha$ and $\beta$.

\begin{lemma}\label{rW9:pseudo}
Let $d$ be $d_{\mathcal G}$, with $\mathcal H=\mathcal G$, or
$d_{\mathcal{GZ}}$, with $\mathcal H$ the group generated by $\mathcal G$ and
$\mathcal Z$.
\begin{enumerate}[label=\textup{(\roman*)},leftmargin=*]
\item $d$ is a lower semicontinuous pseudometric with $d\le\lvert\mathcal
Q\rvert$, and $d(hU,h'U')=d(U,U')$ for all $h,h'\in\mathcal H$.
\item If $d(U,U')=k$, there are $U=U^0,U^1,\dots,U^k$ with $d(U^{i-1},U^i)\le
1$ for $1\le i\le k$ and $d(U^k,U')=0$.
\item $\Gamma_x(hU,\cdot)$ is the image of $\Gamma_x(U,\cdot)$ under $h$ for
$h\in\mathcal H$; hence $W_1^d(\Gamma_x(U,\cdot),\Gamma_x(U',\cdot))=0$
when $d(U,U')=0$.
\item Every bounded real $f\in\mathcal A$ if $d=d_{\mathcal G}$, every bounded
real $f\in\mathcal A_0$ if $d=d_{\mathcal{GZ}}$, has a representative that
is invariant under $\mathcal H$ everywhere and satisfies $\lvert f(U)-f(U')
\rvert\le2\lVert f\rVert_\infty\,d(U,U')$.
\end{enumerate}
\end{lemma}

\begin{proof}
(i) A gauge transformation acts link by link through bijections of $G$, and
so does each $Z^\nu_t$; hence $(hU)_e\ne(hU')_e$ if and only if $U_e\ne
U'_e$, and $\mathrm H_s(hU,hU')=\mathrm H_s(U,U')$. Every $Z\in\mathcal Z$
commutes with every $g\in\mathcal G$, because $z$ is central, so $\mathcal
H=\mathcal G\mathcal Z$ and $d(U,U')=\min_{h\in\mathcal H}\mathrm H_s(U,hU')$.
Then $d(U,U'')\le\mathrm H_s(U,h_1U')+\mathrm H_s(h_1U',h_1h_2U'')=\mathrm
H_s(U,h_1U')+\mathrm H_s(U',h_2U'')$ gives the triangle inequality, and
symmetry and the invariance follow from the same identity. Each indicator
$1[U_Q\ne U'_Q]$ is that of an open set, so $(U,U',h)\mapsto\mathrm H_s(U,
hU')$ is lower semicontinuous; if $(U_k,U'_k)\to(U,U')$ with $d(U_k,U'_k)\le
a$, minimisers $h_k$ have a convergent subsequence in the compact group
$\mathcal H$, with limit $h$, and $\mathrm H_s(U,hU')\le a$.

(ii) Let $U''=hU'$ with $\mathrm H_s(U,U'')=k$, and change the $k$ cells where
$U$ and $U''$ differ one at a time.

(iii) For $g\in\mathcal G$ this is the gauge covariance of
Sec.~\ref{ssec:averaged}, and for $Z^\nu_t$ it is
Proposition~\ref{prop:bpcentre}(i), applied to the kernel. If $d(U,U')=0$,
then $U=hU'$, the image of $\Gamma_x(U',\cdot)$ under $h$ is
$\Gamma_x(U,\cdot)$, and the coupling $(W',hW')$ costs $d(hW',W')=0$.

(iv) Average $f$ over the Haar measure of $\mathcal H$; for each $h$, $f\circ
h=f$ $\mu$-almost everywhere (for $Z^\nu_t$ by
Proposition~\ref{prop:bpcentre}(ii) and the definition of $\mathcal A_0$),
so by Fubini's theorem the average is a representative, and it is invariant
everywhere. If $d(U,U')=0$, then $U=hU'$ and $f(U)=f(U')$; otherwise $d(U,
U')\ge1$ and $\lvert f(U)-f(U')\rvert\le2\lVert f\rVert_\infty$.
\end{proof}

\begin{definition}\label{rW9:def}
Let $n\ge2$ and $s\ge1$, and let $d$ be $d_{\mathcal G}$ or $d_{\mathcal{GZ}}$.
\begin{enumerate}[label=\textup{(\alph*)},leftmargin=*]
\item The \emph{response} $R_d(n,s)$ is the supremum of the distance
$W_1^d(\Gamma_x(U,\cdot),\Gamma_x(U',\cdot))$ over the points $x\in X$, the
cells $Q$ that touch $x$, and the pairs $U,U'$ that agree off $Q$. For a
number $R$ put $\bar R=\max(R,1)$.
\item For $C\ge0$ and $m>0$, the condition $\mathrm{SM}_s(n;C,m)$ holds if for
all $x$, $Q$, $U$, $U'$ as in (a) there are a coupling $\pi$ of
$\Gamma_x(U,\cdot)$ and $\Gamma_x(U',\cdot)$ and a measurable map $h:G^E
\times G^E\to\mathcal G$ such that for every cell $Q'\ne Q$
\[
\pi\bigl\{(W,W'):W_{Q'}\ne(h(W,W')W')_{Q'}\bigr\}\le C\,e^{-ms(\lvert Q'-Q
\rvert_s-1)} .
\]
\item The condition $\mathrm{SM}^{\rm TV}_s(n;C,m)$ holds if for all $x$,
$Q$, $U$, $U'$ as in (a) and every $j\ge1$ the marginals of $\gamma^0_{V_x}
(\cdot\mid U)$ and $\gamma^0_{V_x}(\cdot\mid U')$ on $\Delta_j=V_x\cap
\bigcup\{Q':\lvert Q'-Q\rvert_s\ge j\}$ are at total variation distance at
most $C e^{-ms(j-1)}$.
\end{enumerate}
\end{definition}

The response is the expected number of cells, counted modulo gauge
transformations, in which the two configurations differ after one move,
under the best coupling; $\mathrm{SM}_s$ says that a change of the exterior
in one cell can be absorbed, up to one gauge transformation, everywhere
except with exponentially small probability at grid distance $j$, uniformly in
the exterior. $\mathrm{SM}^{\rm TV}_s$ is the strong mixing condition
of~\cite{MO94a,Mar99} for the raw links, with single-cell changes of the
boundary condition, restricted to the nested sets $\Delta_j$. Since $ms(j-1)$
is a lower bound for the lattice distance between $Q$ and the cells at grid
distance $j$, both conditions express decay at the rate $m$ in lattice
units. Put, for $\lambda>0$ and $q=e^{-\lambda}$,
\[
\begin{split}
\Phi(\lambda)&=\sum_{j\ge1}(64j^3+16j)\,e^{-\lambda(j-1)}\\
&=\frac{64(1+4q+q^2)}{(1-q)^4}+\frac{16}{(1-q)^2},
\end{split}
\]
a decreasing function with $\Phi(\lambda)\to80$ as $\lambda\to\infty$; the
closed form follows from $\sum_{j\ge1}j^3q^{j-1}=(1+4q+q^2)/(1-q)^4$ and
$\sum_{j\ge1}jq^{j-1}=(1-q)^{-2}$.
The three forms are related as follows. The passage from total variation
bounds on nested marginals to one coupling uses a construction of
Goldstein's type~\cite{Gol79}, in the following elementary form with
constant two.

\begin{lemma}\label{rW9:nested}
Let $V$ be a finite set, $V=D_0\supseteq D_1\supseteq\dots\supseteq D_J
\supseteq D_{J+1}=\emptyset$, and let $\alpha,\beta$ be probability measures
on $G^V$ with continuous densities with respect to the Haar measure. Let
$t_j$ be the total variation distance of their marginals on $G^{D_j}$. There
is a coupling of $\alpha$ and $\beta$ under which $\sigma_{D_j}\ne
\sigma'_{D_j}$ has probability at most $2\sum_{i=j}^Jt_i$, for $0\le j\le J$.
\end{lemma}

\begin{proof}
\emph{One step.} Let $\Omega_1=G^{A}$, $\Omega_2=G^{B}$ with $A\cap B=
\emptyset$, $\lambda_1,\lambda_2$ the Haar measures, and $\mu_1,\nu_1$
probability measures on $\Omega_1\times\Omega_2$ with jointly measurable
densities $F,H$. Let $a,b$ be the densities of their marginals on $\Omega_1$,
and $\alpha_y=F(y,\cdot)/a(y)$, $\beta_y=H(y,\cdot)/b(y)$ where these are
defined (an arbitrary fixed density elsewhere). For densities $p,q$ on
$\Omega_2$ let $M(p,q)$ be the maximal coupling: with $w=\int p\wedge q\,
d\lambda_2$, it is the image of $(p\wedge q)\lambda_2$ under $z\mapsto(z,z)$
plus, if $w<1$, $(1-w)^{-1}(p-q)^+\lambda_2\otimes(q-p)^+\lambda_2$; its
marginals are $p\lambda_2$ and $q\lambda_2$, it gives $z\ne z'$ probability
$\int(p-q)^+d\lambda_2$, and it depends measurably on $(p,q)$. Given a
coupling $\pi_1$ of $a\lambda_1$ and $b\lambda_1$, sample $(y,y')$ from
$\pi_1$, then $(z,z')$ from $M(\alpha_y,\beta_y)$ if $y=y'$ and from
$\alpha_y\lambda_2\otimes\beta_{y'}\lambda_2$ otherwise. In both cases $z$
has the law $\alpha_y\lambda_2$ given $y$, and $z'$ the law
$\beta_{y'}\lambda_2$ given $y'$, so this is a coupling of $\mu_1$ and $\nu_1$.
Let $D$ be the image under $(y,y')\mapsto y$ of the restriction of $\pi_1$ to
$\{y=y'\}$. Then $D(S)\le\int_Sa\,d\lambda_1$ and $D(S)\le\int_Sb\,d
\lambda_1$ for every Borel $S$, so $D$ has a density at most $a\wedge b$, and
\[
\begin{split}
\Pr\bigl((y,z)\ne(y',z')\bigr)&\le\pi_1(y\ne y')\\
&\quad+\int(a\wedge b)(y)\int(\alpha_y-\beta_y)^+d\lambda_2\,
d\lambda_1(y).
\end{split}
\]
Pointwise, $(a\wedge b)(\alpha_y-\beta_y)^+\le(F-H)^++(b-a)^+\beta_y$: if
$a\le b$, $a\alpha_y-a\beta_y=(F-H)+(b-a)\beta_y$, and if $a>b$, $b\alpha_y-b
\beta_y\le F-H$. Integrating, the double integral is at most the total
variation distance of $\mu_1,\nu_1$ plus that of their marginals on
$\Omega_1$.

\emph{Induction.} Apply the step with $A=D_{j+1}$, $B=D_j\setminus D_{j+1}$
and the marginals of $\alpha,\beta$ on $G^{D_j}$, for $j=J,J-1,\dots,0$,
starting from the trivial coupling on $G^{D_{J+1}}$. With $p_j$ the
probability of a disagreement on $D_j$, $p_j\le p_{j+1}+t_j+t_{j+1}$ and
$p_{J+1}=t_{J+1}=0$, so $p_j\le2\sum_{i=j}^Jt_i$.
\end{proof}

\begin{proposition}\label{rW9:hierarchy}
Let $n\ge2$, $s\ge1$, $C\ge0$ and $m>0$.
\begin{enumerate}[label=\textup{(\roman*)},leftmargin=*]
\item $\mathrm{SM}^{\rm TV}_s(n;C,m)$ implies $\mathrm{SM}_s(n;2C/(1-e^{-ms}),
m)$, with $h$ the identity.
\item $\mathrm{SM}_s(n;C,m)$ implies
\[
R_{d_{\mathcal{GZ}}}(n,s)\le R_{d_{\mathcal G}}(n,s)\le1+C\,\Phi(ms).
\]
\item Under $\mathrm{SM}_s(n;C,m)$, for all $x$, $Q$, $U$, $U'$ as in
Definition~\ref{rW9:def}\textup{(a)} and every bounded $F$, invariant under
$\mathcal G$ everywhere, that depends only on the links of a set $\mathcal S$
of cells with $Q\notin\mathcal S$,
\[
\bigl\lvert\mathcal P_xF(U)-\mathcal P_xF(U')\bigr\rvert\le2\lVert F\rVert_
\infty\sum_{Q'\in\mathcal S}C\,e^{-ms(\lvert Q'-Q\rvert_s-1)} .
\]
\end{enumerate}
\end{proposition}

\begin{proof}
(i) The configurations of $\Gamma_x(U,\cdot)$ and $\Gamma_x(U',\cdot)$ are
$U$ and $U'$ off $V_x$, which agree off $Q$. Couple the samples of the two
cube laws by Lemma~\ref{rW9:nested} with $D_0=V_x$ and $D_j=\Delta_j$ for
$j\ge1$, and keep
the exteriors. A cell $Q'\ne Q$ at grid distance $j\ge1$ can differ only in
$Q'\cap V_x\subset\Delta_j$, which happens with probability at most
$2\sum_{i\ge j}Ce^{-ms(i-1)}=2Ce^{-ms(j-1)}/(1-e^{-ms})$.

(ii) $d_{\mathcal{GZ}}\le d_{\mathcal G}$, and under the coupling of (b),
$d_{\mathcal G}(W,W')\le\mathrm H_s(W,hW')$, which is at most one for the
cell $Q$ plus the number of cells $Q'\ne Q$ with $W_{Q'}\ne(hW')_{Q'}$.
Taking expectations and counting the cells at grid distance $j$ gives
$1+\sum_{j\ge1}(64j^3+16j)Ce^{-ms(j-1)}$.

(iii) Here $F$ is taken invariant under $\mathcal G$ everywhere, as
is the representative of Lemma~\ref{rW9:pseudo}(iv); that representative is
an average of $F\circ g$ over $g\in\mathcal G$ and so depends only on the
links of the cells of $\mathcal S$. Under the coupling, $F(W')=F(hW')$, and
$F(hW')=F(W)$ unless $W$ and $hW'$ differ on a cell of $\mathcal S$. So
$\lvert\mathcal P_xF(U)-\mathcal P_xF(U')\rvert=\lvert\E_\pi(F(W)-F(hW'))
\rvert$ is at most $2\lVert F\rVert_\infty$ times the probability of such a
difference.
\end{proof}

Part (iii) is the form of the condition stated for gauge-invariant local
functions: a change of the exterior in one cell moves the conditional mean of
every bounded gauge-invariant function of cells at grid distance at least $j$
by at most a multiple of $e^{-ms(j-1)}$, uniformly in the exterior. The
coupling form (b) asks more: one coupling and one gauge transformation for
all cells at once. Whether the gauge-invariant bounds of (iii) imply (b) is
not decided here; patching gauge transformations chosen region by region
creates disagreements along the seams.

\begin{theorem}[strong mixing at one scale implies the averaged
condition]\label{rW9:main}
Let $2\le n$ and $1\le s\le n$, on $\Lambda_L$ with $s\mid L$ or on a box, and
put, with $\bar R=\max(R,1)$,
\[
\varpi(R,n,s)=\frac{\bar R\,(n-s+1)^4-(\bar R-1)(n+s+1)^4}{n^4}.
\]
\begin{enumerate}[label=\textup{(\roman*)},leftmargin=*]
\item If $\varpi=\varpi(R_{d_{\mathcal G}}(n,s),n,s)>0$, then
$\mathrm{BP}(n,\varpi)$ holds.
\item Let $G=\SU(N)$ on $\Lambda_L$. If $\varpi_0=\varpi(R_{d_{\mathcal{GZ}}}
(n,s),n,s)>0$, then $\mathrm{BP}_0(n,\varpi_0)$ holds, and hence
$\mathrm{BP}_{\rm loc}(n,\varpi_0)$.
\end{enumerate}
In both cases, with $\theta=s/n$,
\[
\varpi(R,n,s)\ge(1-\theta)^4-8\theta\Bigl(1+\frac1n\Bigr)\Bigl(\Bigl(1+
\frac1n\Bigr)^2+\theta^2\Bigr)(\bar R-1),
\]
and under $\mathrm{SM}_s(n;C,m)$ one may take $\bar R-1\le C\,\Phi(ms)$.
\end{theorem}
\begin{proof}
(i) Write $d=d_{\mathcal G}$, $R=R_d(n,s)$ and $\kappa_0=\varpi n^4/\lvert X
\rvert$, so that $\kappa_0\lvert X\rvert=\bar R(n-s+1)^4-(\bar R-1)(n+s+1)^4$.

\emph{Adjacent pairs.} Let $d(U,U')\le1$. If $d(U,U')=0$, couple $\mathcal K(U,
\cdot)$ and $\mathcal K(U',\cdot)$ by choosing the same $x$ and using
Lemma~\ref{rW9:pseudo}(iii); the cost is $0$. If $d(U,U')=1$, take $g\in
\mathcal G$ with $\mathrm H_s(U,gU')=1$ and put $U''=gU'$; $U$ and $U''$ differ
exactly on the links of one cell $Q=Q_c$. By Lemma~\ref{rW9:pseudo}(iii), a
coupling $\pi$ of $\Gamma_x(U,\cdot)$ and $\Gamma_x(U'',\cdot)$ gives the
coupling $(W,g^{-1}W'')$ of $\Gamma_x(U,\cdot)$ and $\Gamma_x(U',\cdot)$ with
the same cost, since $d(W,g^{-1}W'')=d(W,W'')$. Choose the same $x$ for both
chains, uniformly in $X$, and couple the two moves at $x$ as follows. If $Q
\subset V_x$, the two configurations agree off $V_x$, the kernels coincide,
and the diagonal coupling costs $0$. If $Q\cap(V_x\cup\partial V_x)=
\emptyset$, the two cube laws coincide, since $\gamma^0_{V_x}(\cdot\mid U)$
depends only on $U_{\partial V_x}$; with one common sample the new
configurations differ only on $Q$, at cost at most $1$. If $x$ and $Q$ touch,
take a coupling of cost at most $R+\eta$, $\eta>0$. Let $N_{\rm in}$ and
$N_{\rm t}$ be the numbers of $x$ in the first and in the third case. The
expected cost is at most
\[
1-\frac{N_{\rm in}-N_{\rm t}(R+\eta-1)}{\lvert X\rvert}.
\]
If $x\in c+[s-n,0]^4$, then $c+[0,s)^4\subset x+[0,n)^4$ and $Q\subset
V_x$; these $(n-s+1)^4$ points are distinct modulo $L$, and they lie in $X$
because $V_x\supset Q\ne\emptyset$, so $N_{\rm in}\ge(n-s+1)^4$. A link of
$V_x\cup\partial V_x$ has its base point in $x+[-1,n]^4$, so if $Q$ meets
$V_x\cup\partial V_x$ then $x\in c+[-n,s]^4$, and $N_{\rm in}+N_{\rm t}\le
(n+s+1)^4$. If $R\ge1$, then $N_{\rm in}-N_{\rm t}(R-1)\ge RN_{\rm in}-(R-1)
(n+s+1)^4\ge\kappa_0\lvert X\rvert$; if $R<1$, then $N_{\rm in}-N_{\rm t}(R-1)
\ge N_{\rm in}\ge\kappa_0\lvert X\rvert$. Letting $\eta\to0$, for $d(U,U')\le1$,
\[
W_1^d\bigl(\mathcal K(U,\cdot),\mathcal K(U',\cdot)\bigr)\le(1-\kappa_0)\,
d(U,U').
\]

\emph{All pairs.} $d$ is a Borel pseudometric (Lemma~\ref{rW9:pseudo}(i)), so
$W_1^d$ satisfies the triangle inequality, by the gluing
lemma~\cite[Ch.~1]{Vil09}. For $d(U,U')=k$, the chain of
Lemma~\ref{rW9:pseudo}(ii) gives $W_1^d(\mathcal K(U,\cdot),\mathcal K(U',
\cdot))\le(1-\kappa_0)k$.

\emph{The spectral step.} Let $f\in\mathcal A$ be real and bounded, with the
representative of Lemma~\ref{rW9:pseudo}(iv), and $\mathrm{Lip}(f)\le2\lVert
f\rVert_\infty$ its Lipschitz constant for $d$. Then $\mathcal Kf$ is bounded
and invariant under $\mathcal G$, by Lemma~\ref{rW9:pseudo}(iii), and for
every coupling $\pi$ of $\mathcal K(U,\cdot)$ and $\mathcal K(U',\cdot)$,
$\lvert\mathcal Kf(U)-\mathcal Kf(U')\rvert\le\mathrm{Lip}(f)\int d\,d\pi$;
so $\mathrm{Lip}(\mathcal K^tf)\le(1-\kappa_0)^t\mathrm{Lip}(f)$ and, as $d\le
\lvert\mathcal Q\rvert$, the oscillation of $\mathcal K^tf$ is at most
$2(1-\kappa_0)^t\lVert f\rVert_\infty\lvert\mathcal Q\rvert$. Since $\mu$ is
invariant under every $\mathcal P_x$, $\mu(\mathcal K^tf)=\mu(f)$, and with
$h=f-\mu(f)$ and $\nu_h$ the spectral measure at $h$ of the self-adjoint
operator $0\le\mathcal K\le I$,
\[
\int\lambda^{2t}d\nu_h=\Var_\mu(\mathcal K^tf)\le4(1-\kappa_0)^{2t}\lVert f
\rVert_\infty^2\lvert\mathcal Q\rvert^2 .
\]
So $\nu_h((1-\kappa_0,1])=0$ and $\ip h{(I-\mathcal K)h}\ge\kappa_0\lVert h
\rVert_2^2$. The truncations $(f\wedge M)\vee(-M)$ of a real $f\in\mathcal A$
lie in $\mathcal A$ and converge to $f$ in $L^2(\mu)$, and both sides are
continuous in $L^2(\mu)$; hence $\ip f{(I-\mathcal K)f}\ge\kappa_0\Var_\mu(f)$
for every real $f\in\mathcal A$; this is the spectral step of
Corollary~\ref{cor:bpcoupling}, where no density argument is needed because
every bounded function of $\mathcal A$ is Lipschitz for $d$. As each
$\mathcal P_x$ is an orthogonal
projection, $\ip f{(I-\mathcal K)f}=\lvert X\rvert^{-1}\sum_x\mu((f-\mathcal
P_xf)^2)$, which is $\mathrm{BP}(n,\kappa_0\lvert X\rvert/n^4)=\mathrm{BP}(n,
\varpi)$.

(ii) The same proof with $d_{\mathcal{GZ}}$, $\mathcal H=\mathcal G\mathcal Z$
and $\mathcal A_0$ in place of $d_{\mathcal G}$, $\mathcal G$ and $\mathcal A$:
the first step takes $h\in\mathcal H$ with $\mathrm H_s(U,hU')=1$, $\mathcal K$
commutes with $\mathcal H$ by Lemma~\ref{rW9:pseudo}(iii), bounded real
functions of $\mathcal A_0$ are Lipschitz by Lemma~\ref{rW9:pseudo}(iv), and
truncation preserves $\mathcal A_0$. This gives $\mathrm{BP}_0(n,\varpi_0)$,
and Theorem~\ref{thm:bplocal}(i) gives $\mathrm{BP}_{\rm loc}(n,\varpi_0)$.

\emph{The bound.} Since $s\le n$, $(n-s+1)^4\ge n^4(1-\theta)^4$, and
$(n+s+1)^4-(n-s+1)^4=8s(n+1)^3+8s^3(n+1)$. The last claim is
Proposition~\ref{rW9:hierarchy}(ii).
\end{proof}

At $\beta=0$ the cube laws do not depend on the exterior, the diagonal
coupling gives $R_{d_{\mathcal G}}(n,1)\le1$, and the theorem returns
$\mathrm{BP}(n,1)$ on $\mathcal A$, as in Proposition~\ref{prop:avgsup}(ii).
The theorem uses the condition only on the cubes of side $n$,
and the cell size enters only through the count of the cubes that contain or
touch a cell; this is what makes the constant physical.

\begin{corollary}[the physical scale]\label{rW9:physical}
Let $\bar R=\max(R_{d_{\mathcal G}}(n,s),1)$.
\begin{enumerate}[label=\textup{(\roman*)},leftmargin=*]
\item If $n\ge8$, $8s\le n$
and $(\bar R-1)s\le n/50$, then $\mathrm{BP}(n,\frac13)$ holds. For
$G=\SU(N)$ on a torus, the same with $R_{d_{\mathcal{GZ}}}$ in place of
$R_{d_{\mathcal G}}$ gives $\mathrm{BP}_0(n,\frac13)$.
\item Let $\ell>0$, $a_k\to0$, and integers $n_k\ge8$, $s_k\ge1$ with $n_k
a_k\le\ell$ and $8s_k\le n_k$. Suppose that for every $k$ there are $C_k,m_k$
with $50\,C_k\Phi(m_ks_k)\,s_k\le n_k$ such that $\mathrm{SM}_{s_k}(n_k;C_k,
m_k)$ holds for the Wilson measure at $\beta_k$ on $\Lambda_L$ for all $L$ in
an infinite set $\mathcal L_k$ of even multiples of $s_k$ with $L\ge2n_k+4$.
Then every continuum limit as in Theorem~\ref{thm:avgtransfer} satisfies
\[
\sigma(H)\subset\{0\}\cup\Bigl[\frac{16}{1875\,e^2\ell},\infty\Bigr),
\]
and the same holds if $G=\SU(N)$ and the condition is assumed only with $h$
ranging over $\mathcal G\mathcal Z$.
\end{enumerate}
\end{corollary}

\begin{proof}
(i) With $\theta=s/n\le\frac18$ and $1+\frac1n\le\frac98$, the bound of
Theorem~\ref{rW9:main} gives $\varpi\ge(\frac78)^4-8\cdot\frac98\cdot(\frac{81}
{64}+\frac1{64})(\bar R-1)\theta\ge\frac{2401}{4096}-\frac{369}{32}\cdot\frac1
{50}>0.355$.

(ii) By Proposition~\ref{rW9:hierarchy}(ii), $(\bar R-1)s_k\le C_k\Phi(m_k
s_k)s_k\le n_k/50$, so (i) gives $\mathrm{BP}(n_k,\frac13)$ on every $\Lambda_
L$, $L\in\mathcal L_k$, and Theorem~\ref{thm:avgtransfer} applies with
$\varpi=\frac13$. In the second case (i) gives $\mathrm{BP}_0(n_k,\frac13)$,
hence $\mathrm{BP}_{\rm loc}(n_k,\frac13)$, and
Theorem~\ref{thm:bplocal}(iii) applies. The coupling $h$ with values in
$\mathcal G\mathcal Z$ bounds $d_{\mathcal{GZ}}$ as in
Proposition~\ref{rW9:hierarchy}(ii).
\end{proof}

In physical units: if the condition holds with a fixed constant $C$, a
physical rate $m_k=m_{\rm ph} a_k$ and cells of side $s_k=\lceil\ell_c/a_k\rceil$,
and $n_k=\lfloor\ell/a_k\rfloor$, then $m_ks_k\ge m_{\rm ph}\ell_c$, $\Phi(m_ks_k)\le
\Phi(m_{\rm ph}\ell_c)$, and the hypotheses of (ii) hold once
$\ell\ge K(\ell_c+2a_k)$ with $K=\max(8,50\,C\Phi(m_{\rm ph}\ell_c))$, since then
$n_k\ge\ell/a_k-1\ge K(\ell_c/a_k+1)\ge Ks_k$: a cube larger than a
fixed multiple of the cell size. For $\ell_c\ge1/m_{\rm ph}$,
$\Phi(m_{\rm ph}\ell_c)\le\Phi(1)<1086$. No constant depends on $a$. The proof uses
the condition in the following form: (a) uniform in the exterior, since the
contraction is used at every pair of configurations that differ in one cell,
the atypical ones included; (b) of coupling type, with exact agreement of
cells modulo one gauge transformation, so that every bounded gauge-invariant
function is Lipschitz with a constant independent of the volume; (c) posed
only on cubes of side $n$; (d) with the expected number of disagreeing cells
bounded, and not merely the influence on each cell small.

\emph{Where the condition holds and where it fails.} For a pseudometric or
metric $d$ write $R_d$ as in Definition~\ref{rW9:def}(a); $R_{\mathrm H_s}
(n,s)$ is the \emph{raw response}, with couplings of the links themselves and
no gauge transformation, and $R_{d_{\mathcal{GZ}}}\le R_{d_{\mathcal G}}\le
R_{\mathrm H_s}$. The proof of Proposition~\ref{rW9:hierarchy}(ii) with $h$
the identity shows that $\mathrm{SM}^{\rm TV}_s(n;C,m)$ gives $R_{\mathrm H_s}
(n,s)\le1+2C\Phi(ms)/(1-e^{-ms})$.

\begin{proposition}[strong coupling]\label{rW9:strong}
Let
\[
\omega_\varphi=\log\frac{\max\varphi}{\min\varphi},\qquad
c=\tanh\frac{\omega_\varphi}2,
\]
and suppose $18c<1$. Then for every $n\ge2$, on every torus and every
box of Sec.~\ref{ssec:averaged},
\[
R_{\mathrm H_1}(n,1)\le1+\frac{72c}{1-18c},
\]
and $\mathrm{BP}(n,\varpi)$ holds with $\varpi\ge1-\frac{2160\,c}{(1-18c)\,n}$
whenever the right side is positive. For the Wilson weight of $G\subset\U(N)$,
$c\le\tanh(N^2\lvert\beta\rvert)$.
\end{proposition}

\begin{proof}
\emph{One link.} Let $e,e'$ be partners and $U,U'$ differ only at $e'$. The
conditional laws $\nu,\nu'$ of $U_e$ have densities proportional to $\prod_{p
\ni e}\varphi(U_p)$ and $\prod_{p\ni e}\varphi(U'_p)$, which differ only in
the factor of the one plaquette $p_0$ that contains $e$ and $e'$. So $\rho=d
\nu/d\nu'$ is a constant times $\varphi(U_{p_0})/\varphi(U'_{p_0})$ and takes
values in an interval $[A,B]$ with $B\le e^{2\omega_\varphi}A$; since $\int
\rho\,d\nu'=1$, $A\le1\le B$. On $[A,B]$ the convex function $(r-1)^+$ lies
below its chord, $(r-1)^+\le(B-1)(r-A)/(B-A)$, so
\[
\begin{split}
\lVert\nu-\nu'\rVert_{\rm TV}&=\int(\rho-1)^+d\nu'\le\frac{(B-1)(1-A)}{B-A}
\\
&=\frac{k+1-kA-A^{-1}}{k-1}\le\frac{(\sqrt k-1)^2}{k-1}
\end{split}
\]
with $k=B/A$, because $kA+A^{-1}\ge2\sqrt k$. The right side is $(\sqrt k-1)/(
\sqrt k+1)\le\tanh(\omega_\varphi/2)=c$. For the Wilson weight
$e^{N\beta\ReTr(\cdot)}$, $\lvert\ReTr\rvert\le N$ on $G$, so
$\omega_\varphi\le2N^2\lvert\beta\rvert$.

\emph{A cube.} Let $V$ be finite, $j\notin V$, and $U,U'$ differ only at $j$;
let $\gamma,\gamma'$ be the two laws $\gamma^0_V$. The couplings of $\gamma$
and $\gamma'$ form a weakly compact set on which $\Psi(\pi)=\sum_{e\in V}\pi(
\sigma_e\ne\sigma'_e)$ is lower semicontinuous, each $\{\sigma_e\ne\sigma'_e\}$
being open; let $\pi^*$ minimise it and $\varrho(e)=\pi^*(\sigma_e\ne
\sigma'_e)$. For $e\in V$, replace $(\sigma_e,\sigma'_e)$ by a sample of the
maximal coupling of the two conditional laws of the link $e$ given the
other links, built from their densities as in the proof of
Lemma~\ref{rW9:nested}, hence measurably; the marginals are invariant, so the
result couples $\gamma$ and $\gamma'$. Changing the partners of $e$ one at a
time and using the one-link bound, the new probability of $\sigma_e\ne
\sigma'_e$ is at most $c\sum_{k\in V}\varrho(k)1[k\text{ a partner of }e]+c\,
1[j\text{ a partner of }e]$, and minimality of $\pi^*$ gives $\varrho(e)$ at
most this. Summing over $e\in V$, as every link has at most eighteen partners,
$(1-18c)\sum_{e\in V}\varrho(e)\le18c$, so there is a coupling under which
the expected number of disagreeing links of $V$ is at most $18c/(1-18c)$; this
is the argument of Lemma~\ref{lem:bpblock} for the discrete metric.

\emph{The response.} A cell of side one consists of the four links with one
base point. If $U,U'$ agree off a cell $Q$, the exteriors of $V_x$ differ in
at most four links, and changing them one at a time and gluing the couplings
gives a coupling of the cube laws with at most $72c/(1-18c)$ disagreeing links
of $V_x$ in expectation. The configurations of $\Gamma_x(U,\cdot)$,
$\Gamma_x(U',\cdot)$ then differ in the cell $Q$ and in at most as many other
cells as there are disagreeing links of $V_x$. Finally $\varpi(R,n,1)=1-(\bar
R-1)((n+2)^4-n^4)/n^4$ and $(n+2)^4-n^4=8n^3+24n^2+32n+16\le30n^3$ for $n\ge
2$, and Theorem~\ref{rW9:main} applies.
\end{proof}

So the hypothesis of Theorem~\ref{rW9:main} holds, in the raw form, at
strong coupling. For $\SU(2)$, with $\beta_W=4\beta$, the range
$18\tanh(\lvert\beta_W\rvert)<1$ is $\lvert\beta_W\rvert<\operatorname{artanh}
(1/18)<0.0557$, smaller than the range $\lvert\beta_W\rvert<\frac29$ of
Corollary~\ref{cor:bpranges}(i), which uses Kantorovich coefficients; the
proposition shows only that the condition is not vacuous. By
Theorem~\ref{rW9:main}, if $\mathrm{BP}(n,\varpi)$ fails for every
$\varpi>0$, then $\varpi(R_{d_{\mathcal G}}(n,s),n,s)\le0$ for every $s$.

\begin{proposition}[two failures]\label{rW9:fails}
\begin{enumerate}[label=\textup{(\roman*)},leftmargin=*]
\item Let $G=\SU(2)$, $\beta\ge1$ and $n\ge2$, and put, as in
Corollary~\ref{cor:bpresponse},
\[
X_n(\beta)=\frac{38400\,n^4(\log\beta+6.11)}\beta .
\]
On $\Lambda_{2n+4}$, with $\bar R=\max(R_{d_{\mathcal G}}(n,1),1)$,
\[
\begin{split}
\bar R-1&\ge\bigl(1-X_n(\beta)\bigr)\frac{n^4}{(n+2)^4-n^4}\\
&\ge\bigl(1-X_n(\beta)\bigr)\frac n{30}.
\end{split}
\]
If $\mathrm{SM}_1(n;C,m)$ holds on $\Lambda_{2n+4}$, then $C\Phi(m)\ge
(1-X_n(\beta))\,n/30$. The raw response $R_{\mathrm H_1}(n,1)$ is the same
on every $\Lambda_L$ with $L\ge2n+4$, and $\max(R_{\mathrm H_1}(n,1),1)-1$
satisfies the same lower bound there. If $\mathrm{SM}^{\rm TV}_1(n;C,m)$
holds on some such torus, then $2C\Phi(m)/(1-e^{-m})\ge(1-X_n(\beta))\,
n/30$.
\item Let $G=\U(1)$ with the Villain weight and $\beta>\beta_0$, as in
Corollary~\ref{cor:avgcoulomb}. There are no $n\ge2$, $1\le s\le n$ and $R$ with
$\varpi(R,n,s)>0$ such that $R_{d_{\mathcal G}}(n,s)\le R$ for the measures
with free boundary of all boxes of side at least $2n$; in particular
$\mathrm{SM}_s(n;C,m)$ fails for some such box whenever $\varpi(1+C\Phi(ms),n,
s)>0$.
\end{enumerate}
\end{proposition}

\begin{proof}
(i) If $\varpi=\varpi(R_{d_{\mathcal G}}(n,1),n,1)>0$, Theorem~\ref{rW9:main}
gives $\mathrm{BP}(n,\varpi)$ on $\Lambda_{2n+4}$, and
Theorem~\ref{thm:bpfixed}(ii) gives $\varpi\le X_n(\beta)$; if $\varpi\le0$
this holds trivially, since $X_n(\beta)>0$. Since $\varpi=1-(\bar R-1)((n+2)^4-n^4)/n^4$, the
bound follows, and $(n+2)^4-n^4\le30n^3$ for $n\ge2$. Under
$\mathrm{SM}_1(n;C,m)$, $\bar R-1\le C\Phi(m)$ by
Proposition~\ref{rW9:hierarchy}(ii). With cells of side one, the cells that
meet $V_x\cup\partial V_x$ have base points in $x+[-1,n]^4$; for $L\ge2n+4$
these, the links of $V_x\cup\partial V_x$ and the plaquettes that meet $V_x$
are those of $\Z^4$, and the links of $\partial V_x$ take all values, as in
the proof of Theorem~\ref{thm:bpstrong}(ii). The cost $\mathrm H_1(W,W')$
counts the cell $Q$ and the cells that meet $V_x$, so the supremum defining
$R_{\mathrm H_1}(n,1)$ does not depend on $L\ge2n+4$, and on
$\Lambda_{2n+4}$, $R_{\mathrm H_1}\ge R_{d_{\mathcal G}}$. The last claim
combines this with the bound $R_{\mathrm H_1}(n,1)\le1+2C\Phi(m)/
(1-e^{-m})$.

(ii) Theorem~\ref{rW9:main} holds on boxes, and would give $\mathrm{BP}(n,
\varpi)$ for all those measures, which Corollary~\ref{cor:avgcoulomb}
excludes. The second claim is Proposition~\ref{rW9:hierarchy}(ii).
\end{proof}

Part (i) is the coupling form of Corollary~\ref{cor:bpresponse}: for $2\le
n\le(\beta/(76800(\log\beta+6.11)))^{1/4}$, where $X_n(\beta)\le\frac12$,
some change of the exterior in one cell outside a cube changes, in the best
coupling and modulo gauge, at least $n/60$ further cells in expectation.
Hence $\mathrm{SM}_1(n;C,m)$ fails on $\Lambda_{2n+4}$, and
$\mathrm{SM}^{\rm TV}_1(n;C,m)$ fails on every $\Lambda_L$ with $L\ge2n+4$,
whenever
\[
\frac{120\,C\,\Phi(m)}{1-e^{-m}}<n\le\Bigl(\frac\beta{76800(\log\beta+
6.11)}\Bigr)^{1/4} .
\]
The obstruction is the charged sector of the small torus
(Corollary~\ref{cor:bpcharged}); it does not bound $R_{d_{\mathcal{GZ}}}$,
whose conclusion is $\mathrm{BP}_0$, which Theorem~\ref{thm:bpfixed} does not
constrain. Part (i) concerns cubes with $n^4\le\beta/(76800(\log\beta+
6.11))$; along the tuning $n_k^4$ grows faster than every power of
$\beta_k$, and part (i) does not apply there. Part (ii) shows, as for (D)
and $\mathrm{BP}$, that a proof of the condition has to use the non-abelian
structure of $G$.

\emph{Mixing at all scales.} If the condition holds on cubes of every large
side with the same constants, the gap follows from the influence bounds of
Proposition~\ref{rW9:hierarchy}(iii) alone, at the mixing rate itself, without
the cube chain.

\begin{theorem}[mixing at all scales gives a gap at the mixing
rate]\label{rW9:allscales}
Let $s\ge1$, $\mathcal L$ an infinite set of even multiples of $s$, and
suppose there are $C\ge0$, $m>0$ and $n_0$ such that $\mathrm{SM}_s(n';C,m)$
holds for the Wilson measure of $\Lambda_L$ at $\beta\ge0$ for every $L\in
\mathcal L$ and every $n'$ with $n_0\le n'\le(L-4)/2$. Then every torus limit
$\omega$ along $\mathcal L$ is invariant under translations and the
hypercubic group and reflection positive, and its transfer matrix has
spectral gap at least $m$.
\end{theorem}

\begin{proof}
The invariance and reflection positivity are as in the proof of
Corollary~\ref{cor:avggap}. Let $F$ be a continuous gauge-invariant function
of finitely many links at positive times, and let $f$, $g$ be the real or
imaginary parts of $\vartheta(\bar F)$ and $\tau_tF$. The links of $g$ have
base points in a box $b+[0,r)^4$, and those of $f$ lie at $\ell^\infty$
distance at least $d$ from it, with $r$ depending only on $F$ and $t-c_F\le d$.
Suppose $d\ge s+3$, put $k=d-2$, $n'=2k+r$ and $x=b-(k,k,k,k)$, and let $L\in
\mathcal L$ with $n_0\le n'\le(L-4)/2$.

The base points of the links of $V_x$ form $b+[-k,k+r)^4$, and those of
$\partial V_x$ lie in the layer $b+([-k-1,k+r]^4\setminus[-k,k+r)^4)$, at
distance at least $k+1<d$ from $b+[0,r)^4$. So $g$ depends only on links of
$V_x$, $f$ only on links off $V_x\cup\partial V_x$, and $\Cov_\mu(f,g)=\mu(f
(\mathcal P_xg-\mu(g)))$, with $\lvert\mathcal P_xg-\mu(g)\rvert$ at most the
oscillation of $\mathcal P_xg$. Let $Q_1,\dots,Q_M$ be the cells that meet
$\partial V_x$; each touches $x$. Given $U,U'$, let $U^i$ equal $U'$ on the
links of $Q_1\cup\dots\cup Q_i$ and $U$ elsewhere. Then $U^{i-1}$ and $U^i$
agree off $Q_i$, and $U^M$ agrees with $U'$ on $\partial V_x$, so $\mathcal
P_xg(U^M)=\mathcal P_xg(U')$. Let $\mathcal S$ be the set of cells that meet
the links of $g$. A cell of $\mathcal S$ and a cell $Q_i$ contain base points
at distance at least $k+1$, so their grid points are at distance at least
$k+2-s$ and $s(\lvert Q_i-Q'\rvert_s-1)\ge k+2-2s=d-2s$; in particular $Q_i
\notin\mathcal S$. By Proposition~\ref{rW9:hierarchy}(iii), applied $M$
times,
\[
\lvert\Cov_\mu(f,g)\rvert\le2C\lVert f\rVert_\infty\lVert g\rVert_\infty\,
M\,\lvert\mathcal S\rvert\,e^{2ms}e^{-md},
\]
with $M\le((n'+2)/s+2)^4$ and $\lvert\mathcal S\rvert\le(r/s+2)^4$. As
$n'=2d+r-4$, this is at most $C_F(1+d)^4e^{-md}$ with $C_F$ independent of $L$
and $d$; for the finitely many smaller $d$, $\lvert\Cov_\mu(f,g)\rvert\le2\lVert
f\rVert_\infty\lVert g\rVert_\infty$. With $d\ge t-c_F$ and the four real
combinations, $\lvert\Cov_{\mu_L}(\vartheta(\bar F),\tau_tF)\rvert\le C'_F
(1+t)^4e^{-mt}$ for all $L\in\mathcal L$ large compared with $t$. The limit
along $\mathcal L$ gives $\lvert\omega(\vartheta(\bar F)\tau_tF)-\lvert
\omega(F)\rvert^2\rvert\le C'_F(1+t)^4e^{-mt}$, and
Proposition~\ref{prop:gapdecay} with every $m'<m$ places the spectrum of the
transfer matrix on $\Omega^\perp$ in $[0,e^{-m}]$.
\end{proof}

\begin{corollary}\label{rW9:allscalescont}
Let $m_{\rm ph}>0$, and suppose that for every $k$ the hypothesis of
Theorem~\ref{rW9:allscales} holds at $\beta_k$ with some $s_k$, $C_k$, $n_{0,k}$,
an infinite set $\mathcal L_k$ and $m_k\ge m_{\rm ph} a_k$, and let $\omega_k$ be a
torus limit along $\mathcal L_k$. Then every continuum limit of $(\omega_k,
O_k)$ satisfies $\sigma(H)\subset\{0\}\cup[m_{\rm ph},\infty)$, with $0$ a simple
eigenvalue.
\end{corollary}

\begin{proof}
As in the proof of Theorem~\ref{thm:avgtransfer}, the proof of
Theorem~\ref{thm:transfer} uses its hypothesis only through the gap of
$\mathbb T_k$ on $\Omega_k^\perp$; by Theorem~\ref{rW9:allscales} that gap
is at least $m_{\rm ph} a_k$, and for $t=\tau/a_k$, $m_{\rm ph} a_kt=m_{\rm ph}\tau$.
\end{proof}

\begin{remark}[what the condition is relative to (D), $\mathrm{BP}$ and the
mass gap]\label{rW9:relations}
Write $\mathrm{SM}_\ell$ for the hypothesis of
Corollary~\ref{rW9:physical}(ii) along the tuning: mixing in the coupling
form of Definition~\ref{rW9:def}(b), on the cubes of physical side $\ell$
only, with cells of physical side $\ell_c$, constants $C$ and $m_{\rm ph}$ fixed, and
$\ell\ge K(\ell_c+2a_k)$ as stated after that corollary. The following is
proved.
\begin{enumerate}[label=\textup{(\alph*)},leftmargin=*]
\item $\mathrm{SM}_\ell\Rightarrow\mathrm{BP}(n_k,\frac13)$ on the tori used
$\Rightarrow\mathrm{BP}_{\rm loc}(n_k,\frac13)\Rightarrow$ a gap of at least
$16/(1875e^2\ell)$ in every nontrivial continuum limit
(Theorem~\ref{rW9:main}, Corollary~\ref{rW9:physical},
Theorems~\ref{thm:bplocal} and~\ref{thm:avgtransfer}). The version modulo
the centre twists gives $\mathrm{BP}_0$ and the same gap.
\item Mixing on cubes of all large sides at the rate $m_{\rm ph} a_k$ gives a gap of
at least $m_{\rm ph}$ in every continuum limit (Corollary~\ref{rW9:allscalescont}),
through the gauge-invariant influence bounds and not through $\mathrm{BP}$.
\item For $\SU(2)$ the coupling form with $s=1$ fails on $\Lambda_{2n+4}$,
and the total variation form on every $\Lambda_L$, $L\ge2n+4$, for the cube
sides $n$ displayed after Proposition~\ref{rW9:fails}; each of the three forms
fails on boxes in the Coulomb phase of the Villain theory
(Proposition~\ref{rW9:fails}); the raw response is bounded at strong
coupling (Proposition~\ref{rW9:strong}).
\end{enumerate}
So $\mathrm{SM}_\ell$ implies $\mathrm{BP}$ and $\mathrm{BP}_{\rm loc}$ along
the tuning, and hence the conclusion of Theorem~\ref{thm:avgtransfer}. The
following is not proved, in either direction. (1) The converse implications: from
$\mathrm{BP}$, $\mathrm{BP}_{\rm loc}$ or the mass gap to $\mathrm{SM}_\ell$.
The first two are statements in $L^2(\mu)$ and the third concerns the
infinite-volume state, while $\mathrm{SM}_\ell$ is a supremum over all
exteriors of a cube, including exteriors of large action; the known converses
for spin systems assume uniformity in the boundary
condition~\cite{SZ92a,Mar99}, and for spin systems the passage from weak
mixing, a bound uniform in the boundary condition but not localised at the
changed boundary site, to strong mixing is proved in~\cite{MOS94} for
squares in two dimensions. (2) An implication between $\mathrm{SM}_\ell$ and
(D) or its ungauged form $D^0_L(n,\alpha)$. (D) is posed in the forest
gauge and $D^0_L$ for the raw links, both through Kantorovich coefficients
normalised by the size of the change of the boundary link; the arguments
here derive no bound on those ratios from a coupling with exact agreement,
and no such coupling from a Kantorovich bound. Both $\mathrm{SM}_\ell$ and
$D^0_L(n,\alpha)$, $\alpha<1$, imply $\mathrm{BP}$ (Theorem~\ref{rW9:main},
Proposition~\ref{prop:avgsup}). (3) That a mixing condition at one scale with
small response implies the same condition at all larger scales, the
analogue of the constructive criterion of~\cite{DS85}; the theorem gives
from one scale only $\mathrm{BP}$ and, through Theorem~\ref{thm:avgdecay},
decay in the measure at the rate $\varpi/(e^2q_nn)$, not uniformly in
exteriors. (4) $\mathrm{SM}_\ell$ along the tuning for any non-abelian group.
The constant of $\mathrm{SM}_\ell$ has to be uniform over exteriors of a cube
whose fields are far from the typical small fluctuations at weak coupling,
and nothing proved here bounds the influence of such exteriors uniformly in
$a$.
\end{remark}

\begin{remark}[the $L^2$ route]\label{rW9:ltwo}
The derivations of~\cite{MO94b,Ces01,BCC02} work in $L^2$: if $P_A,P_B$ are
the conditional expectations given the links outside $A$ and outside $B$,
and $P_{A\cup B}$ that given the links outside $A\cup B$, then for two
orthogonal projections $p,q$ and every vector $v$, $\lVert v-pv\rVert^2+\lVert
v-qv\rVert^2\ge(1-\lVert pq\rVert)\lVert v\rVert^2$, because $\lVert pv\rVert^2+
\lVert qv\rVert^2=\ip vu^2+\ip v{u'}^2$ is at most $\lVert v\rVert^2$ times the
largest eigenvalue $1+\lvert\ip{u}{u'}\rvert$ of the Gram matrix of the unit
vectors $u,u'$ along $pv,qv$, and $\lvert\ip
u{u'}\rvert=\lvert\ip u{pqu'}\rvert\le\lVert pq\rVert$. With $p=P_A-P_{A\cup B}$,
$q=P_B-P_{A\cup B}$ and $v=f-P_{A\cup B}f$ this is $\lVert f-P_{A\cup B}f
\rVert^2\le(1-\varepsilon)^{-1}(\lVert f-P_Af\rVert^2+\lVert f-P_Bf\rVert^2)$
with $\varepsilon=\lVert P_AP_B-P_{A\cup B}\rVert$ on $\mathcal A$; this
norm is at most the essential supremum over the exteriors of $A\cup B$ of the
maximal correlation of the links of $A\setminus B$ and of $B\setminus A$
under the conditional law, so here too the mixing enters uniformly in the
exterior.
The multiscale arguments of~\cite{MO94b,Ces01} iterate this inequality
with mixing on regions of all sizes; on a torus with constants independent
of $L$ that includes regions that wrap the torus. This iteration is not
carried out here; the coupling form above is posed at one scale only.
\end{remark}

\begin{remark}[status]\label{rW9:status}
Proved: a gauge-covariant strong mixing condition on the cubes of side $n$
alone, in coupling form modulo gauge transformations and counted in cells of
side $s$, implies $\mathrm{BP}(n,\varpi)$ on every torus with $s\mid L$ and on
boxes, with $\varpi$ bounded below by an explicit function of the response
and of $s/n$, and its version modulo the centre twists implies
$\mathrm{BP}_0$ (Theorem~\ref{rW9:main}); with exponential decay at a
physical rate and physical cells this gives $\mathrm{BP}(n_k,\frac13)$ and,
in every nontrivial continuum limit, a gap of at least $16/(1875e^2\ell)$
with no constant depending on $a$ (Corollary~\ref{rW9:physical}); the total
variation form implies the coupling form, which bounds the response and the
influence of the exterior on bounded gauge-invariant functions
(Proposition~\ref{rW9:hierarchy}); the raw response is bounded at strong
coupling (Proposition~\ref{rW9:strong}); for $\SU(2)$ at large $\beta$ the
condition with fixed constants fails on cubes of the sides displayed after
Proposition~\ref{rW9:fails}, and it fails on boxes in the Villain Coulomb
phase (Proposition~\ref{rW9:fails}); on all large scales it gives a gap at
the mixing rate (Theorem~\ref{rW9:allscales},
Corollary~\ref{rW9:allscalescont}). The argument is the path coupling
of~\cite{BD97,DSVW04} for block dynamics, carried to gauge orbits and cells.
Not proved: the condition along the tuning, for any non-abelian group; any
converse from $\mathrm{BP}$ or the mass gap to it; any implication between
it and (D); the passage from one scale to all scales; and whether the
gauge-invariant influence bounds imply the coupling form. Neither the
condition nor the mass gap is proved here.
\end{remark}
